% This is the ADASS_template.tex LaTeX file, 24 Oct 2025.
% It is based on the ASP general author template file, but modified to reflect the specific
% requirements of the ADASS proceedings.

% This is the aspauthor.tex LaTeX file
% Copyright 2021, Astronomical Society of the Pacific Conference Series
% Revision:  29 November 2021

% To compile, at the command line positioned at this folder, type:
% pdflatex P813
% bibtex P813
% pdflatex P813
% pdflatex P813

\documentclass[11pt,twoside]{article}

% Do NOT use ANY packages other than asp2023.
\usepackage{asp2023}

\aspSuppressVolSlug
\resetcounters

% References must all use BibTeX entries in a .bib file and be cited with \citet{} or \citep{}.
\bibliographystyle{asp2021}

% Running heads: more than 3 authors -> "Surname1 et al.".
\markboth{Di~Stefano et al.}{Modeling the Sky in VLT/ESPRESSO Spectra}

\begin{document}

\title{Modeling the Sky in VLT/ESPRESSO Spectra}

% Authors separated by commas; the last one preceded by ", and".
\author{Simona~Di~Stefano$^{1,2}$, Giovanni~Lacopo$^1$, Guido~Cupani$^1$, and Valentina~D'Odorico$^1$}
\affil{$^1$INAF -- Osservatorio Astronomico di Trieste, Via Tiepolo 11, I-34143 Trieste, Italy; \email{simona.distefano@phd.units.it}}
\affil{$^2$Dipartimento di Fisica, Sezione di Astronomia, Universit\`a degli Studi di Trieste, Via Tiepolo 11, I-34143 Trieste, Italy}

% This section is for ADS Processing. There must be one line per author. paperauthor has 9 arguments.
\paperauthor{Simona~Di~Stefano}{simona.distefano@phd.units.it}{}{Universit\`a degli Studi di Trieste}{Dipartimento di Fisica}{Trieste}{}{34143}{Italy}
\paperauthor{Giovanni~Lacopo}{giovanni.lacopo@inaf.it}{}{INAF}{Osservatorio Astronomico di Trieste}{Trieste}{}{34143}{Italy}
\paperauthor{Guido~Cupani}{guido.cupani@inaf.it}{}{INAF}{Osservatorio Astronomico di Trieste}{Trieste}{}{34143}{Italy}
\paperauthor{Valentina~D'Odorico}{valentina.dodorico@inaf.it}{}{INAF}{Osservatorio Astronomico di Trieste}{Trieste}{}{34143}{Italy}

% Author index entries (surname first). First author first.
%\aindex{Di~Stefano,~S.}
%\aindex{Lacopo,~G.}
%\aindex{Cupani,~G.}
%\aindex{D'Odorico,~V.}

\begin{abstract}
Accurate sky subtraction is essential for precise flux calibration in high-resolution spectroscopy. For the ESPRESSO spectrograph at the Very Large Telescope (VLT), the official ESO pipeline provides two sky-subtraction strategies: ``pixel-by-pixel'' subtraction and sky modeling with the ``smoothed'' method. Pixel-by-pixel subtraction reduces the final signal-to-noise ratio, while the smoothed method systematically underestimates the observed flux. This ``zero-level problem'' appears as a negative offset in the saturated absorption troughs of quasar spectra. We present an alternative approach, the running-median method, which builds a robust model of the sky continuum, and we test it on a sample of 23 high-resolution, high-redshift quasar spectra. The running-median method models the sky continuum correctly and gives results consistent with the smoothed method. This shows that the negative zero level is driven mainly by the propagation of flux errors when exposures are combined, rather than by the sky model. Refined error-propagation frameworks are therefore needed when processing ESPRESSO data.
\end{abstract}

% Subject index entries (keep commented out).
%\ssindex{instruments!individual!ESPRESSO}
%\ssindex{techniques!sky subtraction}
%\ssindex{observatories!ground-based!VLT}
%\ssindex{spectroscopy!Echelle}

\section{Introduction}
High-redshift quasars are powerful cosmic backlights. The hydrogen and metal absorption lines imprinted on their spectra probe the large-scale structure of the Universe, trace the chemical enrichment of the circumgalactic and intergalactic media, and constrain the processes that drive galaxy formation and evolution. These studies require high-resolution spectra with a high signal-to-noise ratio (S/N). The ESPRESSO spectrograph at the Very Large Telescope (VLT) \citep{pepe+21}, with a resolving power of up to $R\sim130\,000$, is currently one of the most powerful instruments for this work, ahead of ELT/ANDES, which is expected in the 2030s. Fully exploiting these data requires accurate reduction pipelines; in particular, sky subtraction, the removal of atmospheric emission from the science frames to isolate the target signal, must be performed with high precision to ensure reliable flux calibration.

\section{Sky Subtraction in ESPRESSO}
ESPRESSO records the science target with fiber~A and, simultaneously, the sky with fiber~B. The ESO pipeline subtracts the sky from every exposure using this simultaneous sky spectrum, in the following steps:

\begin{enumerate}
    \item \textbf{Throughput calibration:} The throughput of fiber~B relative to fiber~A is measured as a function of wavelength in the 2D extracted (S2D) format, using dedicated calibration exposures of the daylight sky.
    \item \textbf{Spectrum scaling:} The S2D sky spectrum (fiber~B) is scaled by the resulting relative-efficiency curve.
    \item \textbf{Resampling:} The scaled sky spectrum is rebinned onto the wavelength grid of fiber~A by spline interpolation of the cumulative flux as a function of extracted pixel, which conserves flux.
    \item \textbf{Final subtraction:} The scaled and resampled S2D sky spectrum is subtracted from the S2D science spectrum.
\end{enumerate}

\noindent Because the sky emission varies with time, these steps are applied to each science exposure before coaddition.

\section{ESPRESSO Pipeline Methods}
The official ESO ESPRESSO pipeline (v3.3.10) offers two options for sky subtraction:

\begin{itemize}
    \item \textbf{Pixel-by-pixel:} The sky spectrum is subtracted directly from the science spectrum, pixel by pixel, removing both the sky continuum and the sky emission lines. The error on the sky-subtracted frame is the sum in quadrature of the uncertainties of the 2D science frame ($\sigma_{\mathrm{S2DA}}$) and of the 2D sky frame ($\sigma_{\mathrm{S2DB}}$):
    $\sigma_{\mathrm{px\mbox{-}by\mbox{-}px}} = \sqrt{\sigma_{\mathrm{S2DA}}^2 + \sigma_{\mathrm{S2DB}}^2}.$

    \item \textbf{Smoothed:} The sky spectrum is first modeled with a sliding average, and the model is subtracted from the science spectrum. Sky emission lines narrower than the smoothing window are not captured by the model and therefore remain in the sky-subtracted spectrum. The method assumes that the uncertainty of the sky model is negligible, so the error on the sky-subtracted frame depends only on the science exposure: $\sigma_{\mathrm{smoothed}} = \sigma_{\mathrm{S2DA}}$.
\end{itemize}

\noindent Pixel-by-pixel subtraction adds the noise of the sky spectrum and degrades the S/N of faint sources. The smoothed method, on the other hand, has been observed to systematically underestimate the flux, leading to a negative zero level. Mitigating this bias is essential for high-resolution spectroscopy, particularly of faint, low-S/N sources.

\section{Running-Median Sky Subtraction}
Following the same principle as the smoothed method, we re-implemented every step of the sky-subtraction algorithm so that each one could be inspected and controlled. Our procedure (hereafter, the running-median method) models the sky continuum as follows:

\begin{enumerate}
    \item \textbf{Relative throughput calibration:} The 2D extracted sky exposure (S2DB) is corrected for the relative throughput of the two fibers ($\sim1\%$).
    \item \textbf{Removal of spikes and sky lines:} Bright sky emission lines and spikes are masked by sigma clipping.
    \item \textbf{Sky-model subtraction:} The sky continuum is modeled with a running median over a window of 200~pixels ($\simeq200$~km\,s$^{-1}$) and subtracted directly.
    \item \textbf{Equalization, combination, and rebinning:} The flux levels of the 2D extracted exposures are equalized, coadded, and rebinned (step of 1~km\,s$^{-1}$) with Astrocook \citep{cupani+20}.
\end{enumerate}

We tested the running-median method on the EQUALS sample of 23 high-redshift quasars \citep[$z\sim3$--4;][]{berg+25}. The residuals between the sky spectrum (S2DB) and the running-median model (Figure~\ref{P813_f1}) confirm that the model correctly reproduces the sky level, and the method yields a higher S/N than the pixel-by-pixel approach. However, the zero flux level in saturated absorption troughs (e.g., in damped Lyman-$\alpha$ systems and sub-DLAs) shows that the flux is still slightly underestimated, as with the smoothed method.

In the running-median spectra, pixels with negative flux ($f<0$) have systematically smaller errors than pixels with positive flux ($f>0$), whereas for the pixel-by-pixel method the two error distributions almost coincide. Since the exposures are combined and rebinned with an inverse-variance weighted average, negative fluctuations receive larger weights and pull the combined flux below zero.

% Figure file must be named <paper>_f<n> (no subfolder). The extension is omitted so that
% pdflatex picks up P813_f1.png (or .pdf) and the editors' latex run picks up P813_f1.eps.
% CHECK the caption: it must describe what the plot actually shows and make sense in black and white.
\articlefigure[width=0.65\textwidth]{P813_f1}{P813_f1}{Residuals between the sky spectrum (S2DB) and the running-median sky model.}

\section{Conclusions}
The offset is therefore not an artifact of the sky model itself, but a direct consequence of standard error propagation. Rather than resolving the discrepancy, this redefines it as a fundamental statistical challenge that demands advanced noise modeling. Consequently, developing refined error propagation frameworks for ESPRESSO data is a critical next step to unlock the precision required for high-resolution echelle spectroscopy in the ELT/ANDES era.

\bibliography{P813}  % For BibTeX

% if we have space left, we might add a conference photograph here. Leave commented for now.
% \bookpartphoto[width=1.0\textwidth]{foobar.eps}{FooBar Photo (Photo: Any Photographer)}

\end{document}
